\documentclass[11pt]{article}
\usepackage[a4paper,margin=25mm]{geometry}
\usepackage{amsmath,amssymb,amsthm,mathtools,booktabs,hyperref,xurl,microtype}
\setlength{\emergencystretch}{3em}
\hypersetup{hidelinks}
\newtheorem{theorem}{Theorem}
\newtheorem{lemma}[theorem]{Lemma}
\newtheorem{corollary}[theorem]{Corollary}
\newcommand{\Q}{\mathbb Q}\newcommand{\Z}{\mathbb Z}
\newcommand{\Tr}{\operatorname{Tr}}\newcommand{\Nm}{\operatorname{N}}
\title{A common numerator compresses reciprocal products\\
in the ordinary-abc quartic transfer}
\author{\(\Z\Q\) Diophantine Project}
\date{24 November 2026}
\begin{document}
\maketitle

\begin{abstract}
We share the physical numerator point in a finite family of ordinary-abc
integer certificates. The input-scaled point coordinates remain private.
Instead we compress their two reciprocal-product fields: their common
part is a quadratic polynomial in the input, and a difference of squares
expresses their product using one private field. Six shared rational
coordinates replace two private coordinates per input. All residuals stay
quadratic. For the thirteen-predicate corrected Jones source this reduces
695 variables to 675, with 594 quadratic residuals. The retained arithmetic
premises, ordinary rational abc, and ineffective defining shift are unchanged.
Every rational zero has an explicit polynomial lift to the predecessor;
completeness uses a new simultaneous choice of elliptic witnesses.
\end{abstract}

\section{Scope and the tuple theorem}

We use the 56-coordinate ordinary-abc interface of~\cite{paired}, with
the shared reference and fixed dyadic contact translations
of~\cite{reference,offsets}. This interface retains the elliptic group,
formal comparison, trace-index, finite prime-pattern, local preparation,
real recurrence, and pole-realization assertions stated in those sources.
Ordinary rational abc is additional. The present argument proves no new
arithmetic foundation. Its pole equations are unchanged; no general-route
pole replacement is silently imported into this different interface.

The fixed positive odd coefficient $B$ must be a sufficiently large
defining member of that interface. The threshold remains ineffective.
The new system is a restriction followed by algebraic compression, so its
soundness requires no increase of that threshold.

\begin{theorem}\label{thm:tuple}
Under the retained premises, for $k\ge1$ the ordinary-abc interface has
a simultaneous quadratic presentation with
\[
                  46k+14\quad\text{rational coordinates},\qquad
                  41k+9\quad\text{quadratic residuals},
\]
including the $k$ distinguished inputs. Every rational solution has
integral inputs. Every tuple in $(\Z\setminus\{0\})^k$ has a rational
solution. Its sum of squared residuals has degree exactly four.
For $k>3$ this improves the preceding $48k+8$ coordinate count.
\end{theorem}

The integer tuple theorem is enough for the corrected Jones reduction:
the signed-branch theorem~\cite{signed} reconstructs positive witnesses
for all thirteen certified inputs. The source instance is a coefficient,
not an additional unknown.

\begin{corollary}\label{cor:675}
The complete table supplied with this paper has 675 rational unknowns
and 594 quadratic residuals. Its instance projection agrees with the
695-variable ordinary-abc predecessor. Its undecidability interpretation
retains that predecessor's arithmetic premises, ordinary rational abc,
and the externally cited Jones coding theorem. It adds no universality
premise for the historically emended source and asserts no unconditional
undecidability result over $\Q$.
\end{corollary}

\section{The common physical numerator}

Set
\[
 K=\Q(\rho),\quad \rho^2=2,\quad \Delta=11-6\rho,\quad
 E:y^2=x^3-\Delta x,\quad R_0=(1+\rho,2\rho).
\]
Write $f(Q)=1/x(Q)$ and $S_Q=y(Q)/x(Q)$ on the allowed chart.
The fixed-difference witnesses for inputs $n_i\ne0$ use
\[
 A_i=[n_i]R_0,\qquad H=[h]R_0,\qquad
 U_i=[n_i]H,\quad V_i=U_i-A_i.
\]
The contact center depends only on $n_i$ and $A_i$:
\[
 T_i=n_i^2\Tr x(A_i),\quad
 M_i=B+2(n_i^8+T_i^2),\quad a_i=M_i^4+16.
\]
For $2^e\ge k$, the retained contact substitution is
$z_i=z+i/2^e$. The sixteen positions of block $i$ are
\[
 \mathcal D_i=\{\pm(a_i+r)-i/2^e:
                    r\in\{-8,-7,-4,-1,1,4,7,8\}\}.
\]
The existing two-adic trace lemma gives $a_i\in\Z_2$, and distinct
offsets have distinct cosets modulo $\Z_2$. Thus the $\mathcal D_i$
are pairwise disjoint, uniformly in the integer tuple. They are fixed
before the prime selection or the numerator choice.

\begin{lemma}\label{lem:common}
The simultaneous completeness proof may choose a single physical
numerator
\[
                  U=[j]R_0,\qquad V_i=U-A_i
\]
for all blocks, retaining their fixed differences, centers, offsets,
and one common reference $H$.
\end{lemma}
\begin{proof}
We specify the extra order condition, its half-generator issue, and
the local approximation, rather than appealing only to density.
At the contact prime $p_d$ put
\[
 u_d=\Nm(\pi_d-1)/4,\qquad v_d=\Nm(\pi_d+1)/4.
\]
The retained Gaussian residues and all-sign order assertion give
\[
 v_2(u_d)=4,\quad v_2(v_d)=0,\quad
 \gcd(u_d,v_e)=1\ (d,e),\quad
 \gcd(v_d,v_e)=1\ (d\ne e).
\]
The last assertion is needed because the opposite orders will receive
different residues $n_i$. It is the same-sign part of the retained
finite prime supplier, recorded explicitly in~\cite[Section 2]{primes}
and used in the common-numerator CRT of~\cite{common}. It is not inferred
from coprimality of the selected primes alone.

These divided orders annihilate $R_0$, not only $P=2R_0$.
Indeed the halving lemma of~\cite[Section 5]{half}, already used in
the ordinary-abc source~\cite[Section 7]{secants}, applies to the prescribed
Gaussian residue $\pi\equiv9\pmod {16}$. It makes the three abscissa
differences $1+\rho$, $2/(1+\rho)$ and $4$ squares in each reduction.
Thus $R_0\in2E(\mathbb F_p)$ there. Full rational two-torsion then
makes $\#E(\mathbb F_p)/4$ an annihilator of $R_0$. No restriction on
the parity of $n_i$ is introduced.

Before selecting primes, choose allowed reference and numerator
neighborhoods in the cyclic orbit. The finitely many functions attached
to $H,H-R_0,U,U-A_i$ exclude only finitely many orbit points; choose
initial indices divisible by 64 outside that set. Shrink their finitely
many real and fixed-place neighborhoods so all the required chart
functions are bounded. Let $M$ preserve those neighborhoods and include
the order of the compact real component group. Take $64\mid M$.
Include all input supports, contact denominators and differences,
bad chart primes, and preparation data in finite avoidance.

The supplier makes the odd parts of all $u_d,v_d$ prime to $M$.
At seven this follows from its fixed residue $1+3\mathrm i$, whose
adjacent norms are $2$ and $6$ modulo seven; every other fixed odd
prime is handled by its permitted avoidance modulus. Choose the prime
values in the short interval used in~\cite{common}, so that
\[
                 \gcd\bigl(p_d,\prod_e u_e v_e\bigr)=1.
\]
For clarity, with contact values $p_d=u-dD_0^2w^2$, a sufficiently
large slope gives $3u/4<p_d<5u/4$. Hasse's bound makes every relevant
even reduction order less than $3u/2$. If $p_d$ divided another such
order, that even order would be at least $2p_d>3u/2$, a contradiction.
This permits arbitrary additional prime-power precision.

Choose the reference index $h$ by its original congruences and real
return. For each $d\in\mathcal D_i$ impose
\[
 j\equiv j_0\pmod M,\qquad j\equiv0\pmod{u_d},\qquad
 j\equiv n_i\pmod{v_d},\qquad
 j\equiv n_i h\pmod{p_d^{L_d}}.
\]
The opposite orders are pairwise coprime and prime to every primary
order; primary residues are all zero. Their common two-part is
compatible with $64\mid j_0$. The selected prime powers are coprime
to every other modulus. Hence the generalized CRT gives a progression
for $j$. Its step preserves the required real component, and the
retained recurrence chooses $[j]R_0$ in the prescribed real neighborhood.
Finite-place neighborhoods remain preserved by $M$.

At the primary slot $U$ reduces to the identity. At the opposite
slot $U-A_i$ reduces to the identity. Moreover $j-n_i h$ is divisible
by both associated annihilators and by $p_d^{L_d}$. As $L_d$ increases,
the formal group makes $U-[n_i]H$ tend to the identity at both slots.
Consequently $(U,U-A_i)$ approximates
$([n_i]H,[n_i](H-R_0))$ to any required finite precision. The existing
source inequalities for
\[
 C_i=n_i S_U-S_H,\quad I_i=n_i S_{U-A_i}-S_{H-R_0},\quad
 J_i=n_i^2f(H)f(U),\quad K_i'=n_i^2f(H-R_0)f(U-A_i)
\]
are open conditions at these allowed local points. Choose each $L_d$
large enough to retain the error valuations at least three, reciprocal
product valuations at least four, and every strict detector inequality.
Zero errors cause no problem: the requirement that their valuation
exceed a fixed bound is still open. At the other slot the points reduce
to the same $\pm R_0,\pm A_i$ as before; their required unit conditions
are preserved.

At fixed places and real infinity the raw tests stay bounded on the
chosen neighborhoods. The same contact limits hold simultaneously:
$c_i\sim a_i/z_i^2$, $W_i-1=O(c_i)$ and $D_i=O(c_i^2)$.
For the actual retained normalizer these imply
\[
 \alpha_i=c_i^{-1}+O(c_i^2),\qquad
 \tau_i=1+O(c_i),\qquad b_i=7c_i\tau_i^2\sim7c_i.
\]
Thus sufficiently large real $z$ gives $0<b_i<1$. At two, the
unchanged trace lemma gives $a_i\in\Q_2^{\times2}$. A sufficiently
deep denominator makes $c_i/(a_i/z_i^2)$ a square unit and $c_i$
a square of arbitrarily large positive even valuation. The unit
$\tau_i$ then gives exactly $b_i\in7\Q_2^{\times2}$, with the
same required depth as before. At seven, sufficiently positive
$v_7(c_i)$ and bounded raw tests give $v_7(\tau_i)=0$, hence
$v_7(b_i)=1+v_7(c_i)>0$. No special equality of the raw tests is
needed at either place.

Outside the fixed preparation set the contact factorization is
unchanged. Its only positive odd valuations of $u_i=1/c_i$ occur
at the selected contacts, already checked. Negative odd valuations
of $u_i$ at other primes give positive odd valuations of $b_i$ by
the original uniform normalizer lemma, which explicitly allows
arbitrary actual point denominators. At even valuations the square
$\tau_i^2$ preserves parity. This proves the square reduced
denominator of every $b_i$, including any new denominators of the
shared numerator point. The private pole witnesses are supplied
independently by the unchanged interface. All differences
$U-(U-A_i)=A_i$, and hence all height centers and contact lists,
stayed fixed throughout.
\end{proof}

\section{Algebraic compression with quadratic equations}

The physical shared numerator does not make the source's scaled
coordinates $X_i S_U$ and $X_i^2 f(U)$ identical. We leave those
coordinates and the full translated-secant chart unchanged.
The useful shared quantity is
\[
                  J=f(H)f(U)\in K,
             \qquad J_i=X_i^2J.
\]
Let $v=v_H$ be the existing common reference coordinate, and introduce
three common field coordinates $J,E,\mathcal K$, with
\begin{equation}\label{eq:common-coefficients}
                         E=vJ,\qquad \mathcal K=J^2.
\end{equation}
They cost six rational coordinates and four quadratic component rows.

For block $i$, let $\mathcal U_i,\mathcal V_i$ be the old hyperbolic
reciprocal-product coordinates. Their physical interpretation in the
Laurent-packed source is
\[
 \mathcal U_i=J_i+c_i\overline{K_i'},\qquad
 \mathcal V_i=J_i-c_i\overline{K_i'}.
\]
Replace these two fields by one private field $L_i$ and use the
polynomial reconstruction
\begin{equation}\label{eq:lift}
        \mathcal U_i=X_i^2J+L_i,\qquad
        \mathcal V_i=X_i^2J-L_i.
\end{equation}
The already stored coordinates $b_{2,i}=X_i^2$ and $b_{4,i}=X_i^4$
are retained. Therefore the first reciprocal-product row becomes
$u_HF_i=b_{2,i}E$, where $u_H$ is affine in the reference coordinates
and $F_i$ is the unchanged scaled reciprocal abscissa. The second
product row depends only on the difference $2L_i$. Finally,
\begin{equation}\label{eq:difference-squares}
           \mathcal U_i\mathcal V_i=b_{4,i}\mathcal K-L_i^2.
\end{equation}
Each new expression is quadratic in retained or new coordinates.
No input inverse, common contact gauge, or new denominator guard is used.

These are all occurrences of $\mathcal U_i,\mathcal V_i$ in the complete
source: the two components of the first product, the two components
of the second product, and the two detector product components.
The supplied checker audits that complete occurrence list, including
the reciprocal-labelled $i$ input and both added fixed-power inputs.

\begin{lemma}\label{lem:lift}
Every rational zero of the compressed system lifts polynomially to a
zero of the entire predecessor. Conversely, every predecessor zero
satisfying $J_i=X_i^2J$ for a common field $J$ has an extension to the
compressed system. These claims include zero input values and all
allowed tangent, zero-ordinate and torsion-difference fibers.
\end{lemma}
\begin{proof}
Use~\eqref{eq:lift} and~\eqref{eq:common-coefficients}. The first and
second product equations are recovered by substitution.
The equality $b_{2,i}^2=b_{4,i}$, still imposed in every block, recovers
\eqref{eq:difference-squares}. Every other equation is unchanged.
All lifts are polynomial, so none fails on a zero fiber. Conversely
take $L_i=(\mathcal U_i-\mathcal V_i)/2$ and set $E=vJ$, $\mathcal K=J^2$.
The asserted common part gives~\eqref{eq:lift}; substitution proves
all new equations. Division by two is the only division.
\end{proof}

To make the certificate auditable, write a field as $a_0+a_1\rho$.
The four common rows are
\[
 E_0-v_0J_0-2v_1J_1,\quad E_1-v_0J_1-v_1J_0,\quad
 K_0-J_0^2-2J_1^2,\quad K_1-2J_0J_1.
\]
Call these $e_0,e_1,k_0,k_1$. Before the harmless primitive row
divisions, the differences between substituted old and new first
product rows are $2b_{2,i}e_0$ and $-2b_{2,i}e_1$.
For the second product rows the differences are zero.
For detector component $j$ the difference is
\[
            -b_{2,i}^2 k_j+(b_{2,i}^2-b_{4,i})K_j.
\]
These exact identities prove ideal membership in explicitly retained
equations; a numerical sample of elliptic points is not used for the lift.

\section{Counts, source projection, and verification}

The fixed-offset predecessor has $48k+8$ coordinates and $41k+5$
rows. Replacing two private fields by one saves $2k$ rational
coordinates. Adding the three common fields costs six. This gives
$46k+14$ coordinates and $41k+9$ rows. Lemma~\ref{lem:lift} proves
soundness by projection to each original integer certificate.
Lemma~\ref{lem:common} supplies a common physical $U$ for every nonzero
integer tuple, hence a common $J=f(H)f(U)$; the converse lift gives
completeness. This proves Theorem~\ref{thm:tuple}.

For the thirteen-predicate Jones source, start with the complete
695/590 table of~\cite{signed}. Exactly 52 old reciprocal-product
coordinates are replaced by 26 private and 6 common coordinates.
Thus the count is $695-52+26+6=675$; four rows are added, giving 594.
The corrected source's canonical positive integer witnesses supply the
nonzero tuple required by Theorem~\ref{thm:tuple}. Conversely the
polynomial lift recovers every original source row with the same
instance. This proves Corollary~\ref{cor:675}.

The sum of squares has degree exactly four: an unchanged input-square
residual has nonzero quadratic part, and the squared leading quadratic
parts cannot cancel over $\Q$. There is no claim of a bijection with
all predecessor auxiliary tuples. The arithmetic completeness step
reselects elliptic witnesses; the algebraic lift applies to every new
rational zero.

The complete coefficient table, deterministic builder, and independent
checker are in the present directory. Run
\begin{center}\small
\texttt{python -B}\enspace\nolinkurl{undecidability/abc_common_numerator_v103/build_common_numerator.py}\\
\texttt{python -B}\enspace\nolinkurl{undecidability/abc_common_numerator_v103/verify_common_numerator.py}.
\end{center}
The checker imports no builder. It verifies all 78 changed rows by the
displayed polynomial identities, all 512 unchanged rows literally,
all four common equations, every removed-coordinate occurrence, every
input-fourth-power relation, the coefficient ring, full coordinate
support, degree, counts, hashes and metadata. Its finite verification
does not prove the retained arithmetic premises or provide a sufficient
numerical $B$.

\small
\begin{thebibliography}{9}\raggedright
\setlength{\itemsep}{2pt}
\bibitem{paired} \(\Z\Q\) Diophantine Project,
\emph{Paired contact coordinates: an ordinary-abc quartic in 56 variables},
19 November 2026,
\nolinkurl{papers/quartic_existence/paired_contact56_20261119_v81.tex}.
The paired chart and Laurent packing derive from Gabe Conant's notes
with Astra 6.2, audited in v78.
\bibitem{reference} Project, \emph{A shared-reference quartic checkpoint},
23 November 2026,
\nolinkurl{undecidability/papers/shared_quartic_checkpoint.tex}.
\bibitem{offsets} Project, \emph{Fixed dyadic translations for simultaneous
contact certificates}, 23 November 2026,
\nolinkurl{undecidability/papers/abc_fixed_contact_offsets_v93.tex}.
\bibitem{common} Project, \emph{A common numerator point for simultaneous
integer certificates}, 23 November 2026,
\nolinkurl{undecidability/papers/common_numerator_transfers.tex},
Sections 4--6; see also
\nolinkurl{undecidability/papers/common_points_quartic_checkpoint.tex}.
\bibitem{primes} Project, \emph{Arithmetic dependencies of the construction
with 29 variables}, 8 November 2026,
\nolinkurl{papers/arithmetic_dependencies/arithmetic_dependencies_20261108_v8.tex},
Section 2 (the retained Sun finite prime-pattern assertion).
\bibitem{half} Project, \emph{Heights of positive weighted elliptic traces},
14 November 2026,
\nolinkurl{papers/trace_heights/weighted_trace_20261114_v27.tex}, Section 5.
\bibitem{secants} Project, \emph{A quartic in 68 variables from translated
secants}, 17 November 2026,
\nolinkurl{papers/quartic_existence/quartic_abc68_20261117_v65.tex}.
\bibitem{signed} Project, \emph{The signed first-Pell branch},
23 November 2026,
\nolinkurl{undecidability/pell_signed_branch_v101/signed_branch_projection.tex};
complete predecessor
\nolinkurl{undecidability/pell_signed_branch_v101/signed_branch_abc_quadrics.json}.
\end{thebibliography}
\end{document}
